\section{Messaggi chiave}

\subsection{Milestone 1}

\begin{itemize}
    \item La costruzione di un dataset per la predizione della bugginess richiede numerose
    decisioni soggettive; documentarle esplicitamente è fondamentale per la riproducibilità.

    \item L'utilizzo delle FV come AV in assenza di AV esplicite riduce il ricorso a
    Proportion, preservando in parte l'informazione originale del ticket, tenendo sempre conto che si tratti di un'ipotesi.

    \item I code smell sono shiftati di una release per riflettere la loro natura causale:
    i smell introdotti in una release non impattano quella stessa release, ma le successive.

    \item È importante non avvicinarsi temporalmente troppo alle release più recenti del progetto, in quanto possibilmente affette da bug dormienti e/o snoring.
\end{itemize}

\subsection{Milestone 2}

\begin{itemize}
    \item RandomForest domina su AUC e NPofB20 indipendentemente dalla strategia di
    balancing: le probabilità ben calibrate (averaging su 100 alberi) lo rendono il
    classificatore più affidabile per la prioritizzazione delle code review.

    \item La feature selection con soglia $0{,}0$ è un no-op su OpenJPA: tutte le 19
    feature portano informazione sulla bugginess. I metodi wrapper migliorano NaiveBayes
    attenuando la violazione dell'indipendenza condizionale, ma sono neutri per
    RandomForest che gestisce la ridondanza internamente.

    \item Il bilanciamento introduce un trade-off irriducibile tra Recall e Kappa:
    Undersampling massimizza il Recall (meno bug sfuggiti, ma più falsi allarmi);
    SMOTE offre il miglior equilibrio complessivo su RF (Kappa $= 0{,}61$, NPofB20 $= 0{,}88$)
    generando istanze sintetiche per interpolazione invece di duplicare quelle esistenti.

    \item NaiveBayes è strutturalmente limitato dalla violazione dell'assunzione di
    indipendenza condizionale: nessuna strategia di balancing né di feature selection
    ne migliora significativamente le performance perché il bottleneck è la correlazione
    tra feature (LOC\_Touched, Churn, LOC\_Added), non la distribuzione delle classi.

    \item La 10$\times$10 CV è il compromesso standard per dataset con poche release:
    rinuncia alla correttezza temporale (possibile temporal leakage) in favore di una
    stima stabile e confrontabile tra configurazioni. Le metriche tendono a essere
    leggermente ottimistiche rispetto all'uso in produzione.
\end{itemize}

\subsection{Milestone 3}

\begin{itemize}
    \item L'impatto dei code smell è limitato: azzerare NSmells previene solo il 5,60\%
    dei bug predetti in B+, contro il 42\% del paper di riferimento (Falessi). La
    differenza dipende dallo strumento di rilevazione (PMD vs SonarQube), dal contesto
    (open source ad alto turnover vs industriale disciplinato) e dallo sfasamento temporale
    della metrica.

    \item NSmells è ridondante rispetto alle metriche di processo: $\rho = 0{,}78$ con LOC,
    $\rho > 0{,}65$ con Churn, LOC\_Touched e Weighted\_Age. La causa è strutturale:
    il ruleset \texttt{codestyle} di PMD rileva violazioni di stile che crescono linearmente
    con la dimensione del file, già catturata dalle metriche di processo.

    \item Lo studio di ablazione conferma la ridondanza: rimuovere NSmells lascia Kappa e
    AUC invariati ($\Delta \approx 0$), con un lieve miglioramento di Precision ($+0{,}004$).
    Una feature ridondante e collineare può addirittura aggiungere rumore marginale.

    \item Il principale driver di bugginess in Apache OpenJPA è la storia di modifica del
    codice (LOC\_Touched, NR, Churn, Nfix), non la qualità strutturale misurata tramite PMD.
    Metriche di smell architetturali (God Class, Feature Envy, rilevabili con SonarQube)
    potrebbero fornire un segnale più indipendente.
\end{itemize}

\subsection{Milestone 4}

\begin{itemize}
    \item Il filtro \texttt{Nsmells $\geq$ 10} è necessario prima della selezione:
    senza di esso la classe a basso ranking avrebbe near-zero smell, rendendo il
    refactoring triviale e il prompt Copilot privo di contenuto.

    \item La selezione meccanica basata sul nome produce classi agli estremi opposti dello
    spettro qualitativo: BrokerImpl (1576 smell strutturali) e LoginDialog (10 smell
    stilistici). Questo contrasto consente di valutare l'LLM su scenari con difficoltà
    molto diverse.

    \item Il report automatico \texttt{smell\_report.txt} con le violazioni PMD complete
    e lo scheletro del prompt Copilot garantisce la riproducibilità della fase di
    refactoring: le istruzioni passate all'LLM sono tracciabili e non dipendono da
    selezioni soggettive.

    \item I risultati di C\textsubscript{1}--C\textsubscript{4} (step~9) sono in attesa
    di produzione: la valutazione dell'efficacia del refactoring automatico richiede
    la scrittura dei test e l'esecuzione di Copilot sulle due classi.
\end{itemize}
